% Generated draft: requires mathematical and source review.
\documentclass[11pt]{article}
\usepackage[margin=25mm]{geometry}
\usepackage{fontspec}
\setmainfont{Latin Modern Roman}
\usepackage{amsmath,amssymb,mathtools,unicode-math}
\setmathfont{Latin Modern Math}
\usepackage{enumitem,needspace,xurl,hyperref}
\setlength{\parindent}{0pt}
\setlength{\parskip}{5pt}
\setlength{\emergencystretch}{3em}
\newcommand{\N}{\mathbb{N}}
\newcommand{\Z}{\mathbb{Z}}
\newcommand{\Q}{\mathbb{Q}}
\newcommand{\R}{\mathbb{R}}
\newcommand{\C}{\mathbb{C}}
\newcommand{\F}{\mathbb{F}}
\newcommand{\A}{\mathbb{A}}
\newcommand{\PP}{\mathbb{P}}
\newcommand{\E}{\mathbb{E}}
\newcommand{\eps}{\varepsilon}
\newcommand{\dd}{\,\mathrm{d}}
\begin{document}
\section*{UnsolvedMath ID: 6700040}
\textbf{Source problem code:} catalogue entry
\subsection{Definitions and mathematical statement}
Let $(X,g)$ be the given Riemannian manifold. Write $\widehat g$ for the sought metric, distinguishing it from the repeated symbol $g$ in the extracted sentence. The source asks for an area-extremal metric, or alternatively a length-extremal metric, quantitatively close to $g$. It does not define either extremality notion or the allowed scalar-curvature comparison class, so these are unresolved source dependencies.

For a fixed identity-map distance convention, $\lambda$-bi-Lipschitz equivalence would mean
\[
 \lambda^{-1}d_g(x,y)\le d_{\widehat g}(x,y)
       \le\lambda d_g(x,y)\qquad(x,y\in X),\quad\lambda\ge1.
\]
If arbitrary identifying diffeomorphisms are allowed instead, the optimization differs. With the identity convention and a specified extremal class $\mathcal E(X)$, a precise candidate parameter is
\[
 \lambda_*(X,g)=\inf\{\lambda\ge1:
  \exists\widehat g\in\mathcal E(X)
  \text{ satisfying the displayed distance inequalities}\}.
\]
The empty infimum is $+\infty$; a minimizing metric need not be assumed to attain it. Separate choices of $\mathcal E$ give the area and length versions.

The question is whether this minimal distortion can be ``effectively'' evaluated. It particularly asks for an expression in terms of sectional-curvature pinching when $\operatorname{sec}_g>0$. For example, when global bounds $0<K_{\min}\le\operatorname{sec}_g\le K_{\max}<\infty$ exist, a common dimensionless pinching parameter is $K_{\min}/K_{\max}$. The source does not choose this global convention, state compactness or a dimension range, prescribe other allowed dependencies, or say whether effective evaluation means an exact value, an explicit upper bound, or an algorithm. No bound in terms of that ratio is asserted here.
\subsection{Short English statement}
How much distortion is needed to replace a given metric by an area- or length-extremal one, and can that distortion be evaluated from curvature pinching in the positive-sectional-curvature case? The extremality and distortion conventions must be recovered before this becomes a definite optimization problem.

The supplied summary describes a quantitative rigidity/stability direction and reports no specific resolution of the effective evaluation. The OPEN-TRIAGE label and historical 2017 passage are retained with current status marked for review.
\subsection{Sources}
\begin{itemize}
\item Local UnsolvedMath registry, permanent ID 6700040, source code catalogue entry, scalar-curvature Question [?42], proposed by Misha Gromov, 2017. Dataset curation and metadata: Ulam AI and UnsolvedMath Contributors, CC BY 4.0.
\item Misha Gromov, \emph{101 Questions, Problems and Conjectures around Scalar Curvature} (2017), printed tag [?42], occurrence 42, PDF page 48, \url{https://www.ihes.fr/~gromov/wp-content/uploads/2018/08/101-problemsOct1-2017.pdf}; normalized pdf-text extraction, accessed 2026-07-29.
\item No exact later theorem or definition of the extremal classes is supplied. Source-specific provenance: the author-hosted PDF is incomplete and publicly available; current status and statement-reuse terms NEEDS\_REVIEW, separately from the dataset licence.
\end{itemize}
\end{document}
